% !TeX spellcheck = en_us
\documentclass[a4paper,12pt,fleqn]{article}

\usepackage[left=1.50cm, right=1.50cm, top=3cm, bottom=2cm, 
headheight= 2.5cm, headsep= 4pt, footskip=25pt]{geometry}
\usepackage{fancyhdr}
\usepackage{lmodern}
\usepackage[table]{xcolor}
\usepackage{amssymb,amstext,amsmath}
\usepackage{booktabs}
\usepackage{bm}
\usepackage{setspace}
\usepackage{graphicx}
\usepackage[colorlinks]{hyperref}
\setlength{\parindent}{0mm}
\setlength{\parskip}{2mm}
\usepackage{enumitem}
\usepackage{titlesec}
\usepackage{gensymb}
\usepackage{pgfgantt}
\usepackage{lscape}
\usepackage[absolute]{textpos}
\usepackage{comment}
\usepackage{cite}
\usepackage{tikz}
\usepackage{pgfplots}
\usetikzlibrary{shapes.geometric}
\usepackage{array}
\newcolumntype{P}[1]{>{\centering\arraybackslash}p{#1}}
\usepackage{longtable}
\usepackage[font={small}]{caption}
\usepackage[most]{tcolorbox}
\usepackage{tikz}
\usepackage{stmaryrd}
\usepackage{wrapfig}
%\usepackage[none]{hyphenat}
\usepackage{ctable}
\usepackage{ragged2e}
\usepackage[official]{eurosym}
\newcommand*\circled[1]{\tikz[baseline=(char.base)]{
		\node[shape=circle,draw,inner sep=1pt] (char) {#1};}}
\usepackage[symbol]{footmisc}
\renewcommand{\thefootnote}{\fnsymbol{footnote}}
\renewcommand{\thempfootnote}{\fnsymbol{mpfootnote}}
\renewcommand*\footnoterule{}

\usepackage{xcolor}
\definecolor{docColor}{cmyk}{1.00	, 1.00  , 0   , 0.40}  % 0.90\% of black
\color{docColor}

\newtcolorbox{myhlbox}[2][]{colback=white, colframe=docColor, width=1.0\linewidth,
	boxrule=1.0pt, left=0.0pt, right=0.0pt, top=4pt, bottom=2.0pt,
	colbacktitle=docColor!20, coltitle=black, enhanced, attach boxed title to top left={xshift=3mm, yshift=-1mm},
	title=#2, #1}

\hypersetup{
	colorlinks=true,
	citecolor=blue,
	filecolor=magenta,
	linkcolor=blue,
	urlcolor=blue,
	pdftitle={SeminarAnnouncements},
	%bookmarks=true,,
}

\graphicspath{{./figs/}}

\definecolor{dgreen}{rgb}{0.0, 0.5, 0.0}

\allowdisplaybreaks

\newcommand*\mystrut[1]{\vrule width0pt height0pt depth#1\relax}
\renewcommand{\thesection}{\arabic{section}.}
\renewcommand{\thesubsection}{\arabic{section}.\arabic{subsection}}
\renewcommand{\thesubsubsection}{\arabic{section}.\arabic{subsection}.\arabic{subsubsection}}
\titleformat{\section}{\sffamily\large\bfseries}{\thesection}{1.0em}{}
\titleformat{\subsection}{\sffamily\large\bfseries}{\thesubsection}{1.0em}{}
\titleformat{\subsubsection}{\sffamily\normalsize\bfseries}{\thesubsubsection}{1.0em}{}
\titlespacing\section{0pt}{0pt plus 2pt minus 2pt}{0pt plus 2pt minus 2pt}
\renewcommand{\baselinestretch}{1.00}

\newcommand\CircledImage[1]{%
	\tikz\path[fill overzoom image={#1}, draw=docColor, line width=0mm]circle[radius=0.5\linewidth];%
}

\input{format/defs}

\fancypagestyle{titlepage}{
	\fancyhf{}
	\insertLXheader
}

\fancypagestyle{plain}{%
	\fancyhf{}
	\renewcommand{\headrulewidth}{0pt}
	\renewcommand{\footrulewidth}{0pt}
	\insertnormalpage
}

\pagestyle{plain}

\begin{document}
	\sffamily
\thispagestyle{titlepage}
\vspace*{-3em}
\begin{center}
	\huge \textbf{LMS Seminar}
\end{center}
\begin{center}
	\Large 
	\textbf{Mechanical behavior at the nanoscale in the light of synchrotron X-ray diffraction} 
\end{center}
\begin{center}
	\Large
	Thomas Cornelius \\
	{\large Institute for Materials, Microelectronics and Nanoscience of Provence \\ 
	Aix-Marseille University} \\
	\vspace*{1em}
	{\Large
	\textbf{Date and Time}:~September 18, 2025 (2 -- 3 pm)
	
	\textbf{Venue}:~Amphi 104 (Pole Meca) }
\end{center}
\vspace*{-1em}
\begin{myhlbox}{{\large Abstract}}
	Strain is known to affect the chemical and physical properties of materials. For nanostructures the limit of elasticity was found to increase with decreasing structure size reaching the ultimate limit of the material for defect-scarce nano-objects. This elevated yield strength opens new pathways in elastic strain engineering in nanomaterials. In addition, defects generating strong strain gradients, thus modifying inter-atomic spacings, may activate emergent physical phenomena. Besides strain, crystal phase transformations may strongly affect the physical and mechanical properties like inducing plasticity (TRIP effect) instead of brittleness in certain types of steels or even ceramics. To further elucidate the mechanical properties at the nanoscale, in situ techniques allowing for visualizing the evolution of strain, defects, and crystalline phases during nanomechanical testing are being developed.
	X-ray diffraction techniques are non-invasive and thus perfect tools to study the mechanical behavior at the nanoscale in situ. Thanks to recent developments on focusing optics, synchrotron hard X-ray beams are nowadays routinely focused down to the 100-nm scale, thus facilitating the investigation of single nanostructures. In the last two decades, Bragg coherent X-ray diffraction imaging, which is a lensless imaging technique, has been developed yielding spatial resolutions of ~10 nm and displacement field resolutions in the pm-range [2]. On the other hand, Laue microdiffraction is highly sensitive to crystal orientations, defects as well as crystalline phases. It further does not need any sample rotation making it the optimal technique for in situ nanomechanical studies.
	Here, examples on in situ mechanical studies of the elastic and plastic behavior as well as phase transformations at the nanoscale will be presented 
	
\end{myhlbox}

\vspace*{1em}

\begin{myhlbox}{{\large About the speaker}}
	\begin{wrapfigure}{r}{0.30\textwidth}
		\vspace*{0em}
		\CircledImage{images/cornelius}
	\end{wrapfigure}
	Thomas W. Cornelius is CNRS research director at the Institute for Materials, Microelectronics and Nanoscience of Provence (IM2NP) in Marseille. After his studies in physics and his PhD at the Ruprecht-Karls University in Heidelberg (Germany), he performed a postdoctoral fellowship at the GSI Helmholtz Center for heavy ion research in Darmstadt (Germany) followed by a postdoc at the ID01 beamline at the European Synchrotron ESRF in Grenoble before joining CNRS in 2011. He is head of the Mechanics of Nano-objects (MNO) research group at IM2NP and further manages the CNRS funded research network GDR CohereX on science with coherent X-rays at 3rd and 4th generation synchrotron sources. From 2016 to 2021, he was member of the CNRS National Committee section 05 and is current member of the CoNRS of section 07. 
	His research focuses on the nanomechanical properties of nanostructures and the structure-property relationship in functional materials at the nanoscale using synchrotron X-ray diffraction methods such as Bragg coherent X-ray diffraction imaging, scanning X-ray diffraction microscopy, or Laue microdiffraction. He develops original in situ experiments allowing for investigating the strain response and plasticity in functional nanomaterials during the application of mechanical, electrical, or thermal actuation to study the structure-property relationship.
	
\end{myhlbox}
\end{document}